\documentclass[a4paper,twoside]{article}


% --- RUIMTE EN LAYOUT ---
\usepackage{setspace}
\setstretch{1.15}

\usepackage{geometry} %size regulating
\geometry{
 a4paper,
 left=25mm,
 right=25mm,
 top=20mm,
 bottom=20mm,
 }

\usepackage{parskip} % for better spacing between paragraphs


% --- WISKUNDE ENSYMBOLEN ---
\usepackage{amsmath, amssymb, amsfonts, amsthm} % For nicer math format
\usepackage{mathrsfs} % For curly arrr
\usepackage{bm} % For bold symbols eg \bm{\nabla}
\usepackage{esint} % For bolletje in integrals
\usepackage{dsfont}
\usepackage{siunitx} % Voeg deze expliciet toe voor die graden en eenheden!


% --- OVERIGE ---

\usepackage{graphicx} % Required for inserting images
\usepackage[hidelinks]{hyperref} % geen kaders om linken
\usepackage{subcaption}
%\usepackage[dutch]{babel}
\usepackage[T1]{fontenc}
\usepackage[utf8]{inputenc}
\usepackage{lmodern}
\usepackage{multicol}

% --- Tikz packages ---
\usepackage{tikz}
\usetikzlibrary{angles,quotes, babel} % for pic (angle labels)
\usetikzlibrary{arrows.meta,decorations.markings,calc, decorations.pathmorphing, decorations.pathreplacing, positioning, shapes, shadings, patterns}
\usepackage{xcolor}
\usepackage{tikz-3dplot}
\usepackage[siunitx]{circuitikz}

\usepackage{etoolbox} % ifthen
\usepackage[outline]{contour} % glow around text
\usepackage{xfp} % higher precision (16 digits?)
\contourlength{1.1pt}



% --- THEOREM STIJLEN MET EXTRA RUIMTE ---

\theoremstyle{definition}
\newtheorem{postulate}{Postulaat}
\newtheorem{definition}{Definitie}
\newtheorem{derivation}{Afleiding}

\theoremstyle{plain}
\newtheorem{theorem}{Theorem}


% --- EXTRA COMMANDO'S ---

\newcommand{\dbar}{{\mkern+3mu\mathchar'26\mkern-12mu d}}

% --- Titel en Introductie --- 

\title{Interpretatie Corollary 3.6.1}
\author{T. Cools}
\date{Groepen en Representaties, 2025-2026}

\begin{document}
\maketitle

\begin{abstract}
    Ik ben een abstract.
\end{abstract}

De great orthogonality theorem zegt:
\[
\frac{1}{\#G}\sum_{x} D_{ij}^{\alpha}(x) \bar{D}_{kl}^{\beta}(x) = \frac{1}{d_\alpha} \delta_{\alpha\beta} \delta_{ik} \delta_{jl}
\]

In \textbf{Corollary 3.6.1} is 
\[Q = \sum_{x} D^\alpha (x) B D^{\beta \dagger} (x)\]


Volgens Schurs lemma:  \[Q = \delta_{\alpha, \beta} \cdot c_\alpha \cdot \mathds{1}_{d_\alpha}\]





Geval \(\alpha = \beta\) zegt:

\[
Q =  c_\alpha \cdot \mathds{1}_{d_\alpha}
\]

En als we het spoor nemen: 
\[tr(Q) = c_\alpha d_\alpha\]

en 

\[tr(Q) = \sum_{x} tr(D^\alpha(x)BD^{\alpha, \dagger}) \]

Frank zei dat \(tr(ABC) = tr(BCA)\), een even permutatie dus:

\[tr(Q) = \sum_{x} tr(BD^{\alpha, \dagger}(x) D^\alpha(x)) = \sum_{x} tr(B) = \#Gtr(B) \]

want er zijn \(\#G\) elementen.

Combineren we wat we hebben gevonden voor \(tr(Q)\):

\[c_\alpha d_\alpha = \#G tr(B) \]

\[\implies c_\alpha = \frac{tr(B) \# G}{ d_\alpha}\]


Tot hier zijn we al geraakt \(\dots\)

\textbf{Moeilijkheid}: we nemen \(B = |j \rangle \langle l|\)

Dan is:

 \[Q = \sum_{x} D^\alpha (x) |j \rangle \langle l| D^{\alpha \dagger} (x)\]

\[\implies Q_{ik} = \sum_{x} D^\alpha_{ij}(x) \bar{D}^\alpha_{kl}(x)\]

Via orthogonaliteit:
\[
Q_{ik}
 = 
\frac{\#G}{d_\alpha} \delta_{\alpha \alpha} \delta_{ik} \delta_{jl} 
=
 \frac{\#G}{d_\alpha} \delta_{ik} \delta_{jl} 
\]



Nemen we het spoor opnieuw:

\[tr(Q) 
=
 \sum_{i=1}^{d_\alpha}Q_{ii} 
 = \sum_{i=1}^{d_\alpha}\frac{\#G}{d_\alpha} \delta_{ii} \delta_{jl}
 = \sum_{i=1}^{d_\alpha}\frac{\#G}{d_\alpha} \delta_{jl}
 = \#G \delta_{jl}
\]

En we hadden al \(tr(Q) = c_\alpha d_\alpha\), dit gecombineerd geeft:

\[c_\alpha d_\alpha = \#G \delta_{jl} \]






\end{document}